% Full LaTeX report for EduLearn
\documentclass[12pt,a4paper]{report}

% -------------------------------------------------
% Packages
\usepackage[utf8]{inputenc}
\usepackage[T1]{fontenc}
\usepackage[french]{babel}
\usepackage[margin=2.5cm]{geometry}
\usepackage{graphicx}
\usepackage{hyperref}
\usepackage{amsmath}
\usepackage{amssymb}
\usepackage{float}
\usepackage{longtable}
\usepackage{booktabs}
\usepackage{listings}
\usepackage{xcolor}
\usepackage{titlesec}
\usepackage{fancyhdr}
\usepackage{setspace}
\usepackage{chngcntr}
\usepackage{microtype}
\usepackage{tocloft}
\usepackage{tikz}
\usepackage{calligra}
\usepackage{biblatex}
\addbibresource{bibliography.bib}

% Counter configuration
\counterwithout{figure}{chapter}
\counterwithout{table}{chapter}

% Font and colors
\usepackage{fontspec}
\setmainfont{Inter}
\definecolor{fsegsblue}{RGB}{0,85,140}
\definecolor{limeCustom}{HTML}{F6F6C4}
\definecolor{limeShadow}{HTML}{C4D600}

% Hyperref configuration
\hypersetup{
    colorlinks=true,
    linkcolor=black,
    filecolor=magenta,
    urlcolor=blue,
    citecolor=black
}

% Title format for chapters
\titleformat{\chapter}[display]
  {\normalfont\huge\bfseries\centering}
  {\chaptertitlename\ \thechapter}
  {20pt}
  {\Huge}

% Fancy header/footer
\pagestyle{fancy}
\setlength{\headheight}{15pt}
\fancyhf{}
\fancyhead[L]{\leftmark}
\fancyfoot[C]{\thepage}
\renewcommand{\headrulewidth}{0.4pt}
\renewcommand{\footrulewidth}{0.4pt}

% Justification and paragraph spacing
\tolerance=2000
\emergencystretch=10pt
\hyphenpenalty=500
\hbadness=2000
\setlength{\parindent}{0pt}
\setlength{\parskip}{8pt plus 2pt minus 1pt}

% Code listing style
\definecolor{codegreen}{rgb}{0,0.6,0}
\definecolor{codegray}{rgb}{0.5,0.5,0.5}
\definecolor{codepurple}{rgb}{0.58,0,0.82}
\definecolor{backcolour}{rgb}{0.95,0.95,0.92}
\lstdefinestyle{mystyle}{
    backgroundcolor=\color{backcolour},
    commentstyle=\color{codegreen},
    keywordstyle=\color{blue},
    numberstyle=\tiny\color{codegray},
    stringstyle=\color{codepurple},
    basicstyle=\ttfamily\footnotesize,
    breakatwhitespace=false,
    breaklines=true,
    captionpos=b,
    keepspaces=true,
    numbers=left,
    numbersep=5pt,
    showspaces=false,
    showstringspaces=false,
    showtabs=false,
    tabsize=2,
    literate={á}{{\'a}}1 {é}{{\'e}}1 {í}{{\'i}}1 {ó}{{\'o}}1 {ú}{{\'u}}1
             {Á}{{\'A}}1 {É}{{\'E}}1 {Í}{{\'I}}1 {Ó}{{\'O}}1 {Ú}{{\'U}}1
             {à}{{\`a}}1 {è}{{\`e}}1 {ì}{{\`i}}1 {ò}{{\`o}}1 {ù}{{\`u}}1
             {À}{{\`A}}1 {È}{{\'E}}1 {Ì}{{\`I}}1 {Ò}{{\`O}}1 {Ù}{{\`U}}1
             {ä}{{\"a}}1 {ë}{{\"e}}1 {ï}{{\"i}}1 {ö}{{\"o}}1 {ü}{{\"u}}1
             {Ä}{{\"A}}1 {Ë}{{\"E}}1 {Ï}{{\"I}}1 {Ö}{{\"O}}1 {Ü}{{\"U}}1
             {â}{{\^a}}1 {ê}{{\^e}}1 {î}{{\^i}}1 {ô}{{\^o}}1 {û}{{\^u}}1
             {Â}{{\^A}}1 {Ê}{{\^E}}1 {Î}{{\^I}}1 {Ô}{{\^O}}1 {Û}{{\^U}}1
             {œ}{{\oe}}1 {Œ}{{\OE}}1 {æ}{{\ae}}1 {Æ}{{\AE}}1 {ß}{{\ss}}1
             {ç}{{\c c}}1 {Ç}{{\c C}}1}
\lstset{style=mystyle}

% One‑half line spacing
\onehalfspacing

\begin{document}

% -------------------------------------------------
% Page de garde
\definecolor{limeShadow}{HTML}{C4D600}
\definecolor{limeCustom}{HTML}{F6F6C4}
\begin{titlepage}
    \newgeometry{left=3.8cm,right=2cm,top=2cm,bottom=2cm}
    \begin{tikzpicture}[remember picture, overlay]
        \node[anchor=north west, inner sep=0pt] at (current page.north west) {
            \includegraphics[height=\paperheight,width=5cm]{figures/bande_gauche.png}
        };
    \end{tikzpicture}
    \begin{center}
        \sffamily\onehalfspacing
        % En‑tête administratif
        {\small \textbf{RÉPUBLIQUE TUNISIENNE}}\\
        \vspace{0.05cm}
        {\footnotesize \textbf{MINISTÈRE DE L'ENSEIGNEMENT SUPÉRIEUR\\
        ET DE LA RECHERCHE SCIENTIFIQUE}}\\
        \vspace{0.05cm}
        {\small \textbf{UNIVERSITÉ DE SFAX}}\\
        \vspace{0.6cm}
        % Logo
        \includegraphics[width=4.5cm]{figures/FSEGPFE.png}\\
        \vspace{0.8cm}
        % Titre du PFE
        {\Large \textbf{PROJET DE FIN D'ÉTUDES}}\\
        \vspace{0.4cm}
        {\normalsize EN VUE DE L'OBTENTION DU DIPLÔME EN}\\
        \vspace{0.4cm}
        {\Large \textbf{LICENCE EN INFORMATIQUE DE GESTION}}\\
        \vspace{0.6cm}
        % Bande de sujet
        \begin{tikzpicture}
            \useasboundingbox (0,0) rectangle (\linewidth,2.2cm);
            \fill[limeCustom] (-1.3cm,0) rectangle (17.2cm,2.2cm);
            \fill[limeShadow] (-1.3cm,-0.5cm) rectangle (17.2cm,0);
            \node[align=center,text width=15cm] at (7.95cm,1.1cm) {
                \begin{minipage}{15cm}
                    \centering\sffamily\setlength{\baselineskip}{18pt}
                    \textbf{\large UNE APPLICATION SYSTÈME WEB LÉGER DE GESTION DE L'APPRENTISSAGE (LMS)}
                \end{minipage}%
            };
        \end{tikzpicture}
        \vspace{0.8cm}
        % Acteurs du projet
        \begin{spacing}{1.4}
            {\large \textbf{ÉLABORÉ PAR : } Abderrahmen Ktari}\\
            \vspace{0.3cm}
            {\large \textbf{ENCADRANTE ACADÉMIQUE : } Mme. Afef Jmal}\\
            \vspace{0.3cm}
            {\large \textbf{ENCADRANTE PROFESSIONNELLE : } Mme. Afef Marweni}\\
            \vspace{0.3cm}
            {\large \textbf{STAGE EFFECTUÉ À : } MTD Group}
        \end{spacing}
        \vfill
        {\large Année universitaire 2025/2026}
    \end{center}
    \restoregeometry
\end{titlepage}
\clearpage
\pagenumbering{roman}

% -------------------------------------------------
% Dédicaces
\clearpage
\addcontentsline{toc}{chapter}{Dédicaces}
\begin{center}
    \vspace*{1.5cm}
    {\Huge\color{fsegsblue}\textbf{\textit{Dédicaces}}}
    \vspace{1.5cm}
    \large\itshape
    C'est grâce à Dieu et à Sa volonté que ce travail a été réalisé.\\
    Avec l'expression de ma plus profonde reconnaissance, je dédie ce projet de mon stage de fin d'études :
    \vspace{1.2cm}
    \large
    À mes chers parents, Samir et Maweheb, pour leurs efforts qui ont mené à ma réussite.\\
    À mes frères et sœur Ahmed, Brahim et Dorra pour leurs encouragements.\\
    Quels que soient les termes employés, je n'arriverai jamais à exprimer ma gratitude sincère, pour leurs encouragements et leurs conseils précieux tout au long de ce parcours. Je leur souhaite tout le succès et le bonheur.
    \vspace{1.2cm}
    \large
    À mes amis et toutes les personnes qui ont croisé mon chemin et qui, à un moment donné, m'ont soutenu et encouragé à réaliser mes aspirations.\\
    Et sans oublier tous mes enseignants qui m'ont guidé tout au long de mon parcours d'étude.
\end{center}
\clearpage

% -------------------------------------------------
% Remerciements
\clearpage
\addcontentsline{toc}{chapter}{Remerciements}
\begin{center}
    \vspace*{1.5cm}
    {\Huge\color{fsegsblue}\textbf{\textit{Remerciements}}}
    \vspace{1.5cm}
    \large\itshape
    C'est avec un grand honneur que j'exprime mes vifs remerciements et ma profonde gratitude envers tous ceux qui ont contribué à l'élaboration de ce projet de fin d'études.
    \vspace{1cm}
    Je tiens tout d'abord à exprimer mes profonds remerciements à mon encadrante académique, Madame Afef Jmal. Son accompagnement, ses conseils, sa disponibilité et ses encouragements constants ont été le moteur de ce projet.
    \vspace{1cm}
    Mes remerciements s'adressent également à Madame Afef Marweni, pour m'avoir accueilli au sein de la société MTD Group pour réaliser mon stage de projet de fin d'études. Et aussi pour la confiance qu'elle m'a accordée, elle n'a cessé de me faire profiter de ses précieux conseils et remarques.
    \vspace{1cm}
    Je souhaite également exprimer ma gratitude à tous mes enseignants de la Faculté des Sciences Économiques et de Gestion de Sfax, qui ont veillé à la qualité de notre formation.
\end{center}
\clearpage

% -------------------------------------------------
% Table des matières
\renewcommand{\contentsname}{Sommaire}
\tableofcontents
\clearpage
% Listes de figures et de tableaux
\renewcommand{\listfigurename}{Table des Figures}
\listoffigures
\addcontentsline{toc}{chapter}{Table des Figures}
\clearpage
\renewcommand{\listtablename}{Liste des Tableaux}
\listoftables
\addcontentsline{toc}{chapter}{Liste des Tableaux}
\clearpage

% -------------------------------------------------
% Introduction générale
\pagenumbering{arabic}
\clearpage
\chapter*{Introduction générale}
\addcontentsline{toc}{chapter}{Introduction générale}
\markboth{Introduction générale}{Introduction générale}
\clearpage
La transformation digitale des systèmes d'enseignement et de formation professionnelle constitue aujourd'hui un enjeu majeur pour les entreprises et les établissements d'enseignement. Dans ce contexte, les plateformes de gestion de l'apprentissage, communément appelées \textit{Learning Management Systems} (LMS), jouent un rôle central en offrant des solutions intégrées pour la diffusion des contenus pédagogiques, l'évaluation des connaissances et le suivi personnalisé des apprenants. Ces outils permettent non seulement de centraliser les ressources éducatives, mais aussi d'automatiser les processus administratifs liés à la formation, améliorant ainsi l'efficacité pédagogique et l'expérience utilisateur.

C'est dans cette dynamique que s'inscrit notre projet de fin d'études, réalisé au sein du groupe MTD Group, dont l'objectif principal est de concevoir et de développer un système web léger de gestion de l'apprentissage. Actuellement, la gestion des formations souffre souvent d'une dispersion des ressources pédagogiques, d'une charge administrative importante liée à la correction manuelle des évaluations et d'une absence d'outils de suivi de la progression des apprenants. Face à ce constat, notre projet vise à proposer une solution centralisée, intuitive et performante, adaptée aux trois profils d'utilisateurs que sont l'administrateur, l'enseignant et l'étudiant.

Le présent mémoire est structuré en quatre chapitres complémentaires, reflétant les différentes phases du cycle de développement de notre projet :
\begin{itemize}
    \item Le premier chapitre, « Définition du Projet et Analyse des Solutions Existantes », présente le cadre du projet, l'organisme d'accueil, les objectifs visés, ainsi qu'une étude comparative des plateformes LMS existantes afin de situer notre contribution.
    \item Le deuxième chapitre, « Capture des besoins », est consacré à la formalisation des exigences fonctionnelles de l'application à l'aide du langage UML, en identifiant les acteurs clés et en détaillant leurs cas d'utilisation respectifs.
    \item Le troisième chapitre, « Analyse et conception », aborde la modélisation technique du système à travers les diagrammes de séquence et de classes, l'architecture logicielle MVC adoptée et la dérivation du schéma physique de la base de données.
    \item Enfin, le quatrième chapitre, « Implémentation », détaille l'environnement de réalisation, les technologies retenues (Laravel, MySQL, Bootstrap) et présente les principales interfaces développées pour les différents utilisateurs de la plateforme.
\end{itemize}

L'ensemble de ce travail démontre la faisabilité d'un LMS léger, intuitif et performant, parfaitement adapté aux besoins de formation. Il constitue une contribution concrète à la transformation digitale des pratiques pédagogiques et une application pratique des compétences acquises durant notre formation.
\clearpage

% -------------------------------------------------
% Chapitre 1 – Définition du Projet et Analyse des Solutions Existantes
\chapter{Définition du Projet et Analyse des Solutions Existantes}

\section{Introduction}
Ce premier chapitre a pour objectif de poser les bases de notre projet de fin d'études. Nous allons définir le cadre général dans lequel s'inscrit notre travail, présenter l'entreprise qui nous accueille, et analyser les solutions déjà disponibles sur le marché des plateformes d'apprentissage en ligne. Nous commencerons par décrire la mission qui nous a été confiée, en expliquant le contexte professionnel et les objectifs à atteindre. Ensuite, nous étudierons le contexte du sujet en identifiant les difficultés liées à la gestion des formations. Puis, nous comparerons des applications existantes comme Moodle et Google Classroom, pour voir leurs points forts et leurs points faibles. Enfin, nous présenterons la problématique à laquelle notre solution, EduLearn, cherche à répondre, avant de conclure ce chapitre et d'annoncer la suite du rapport.

\section{Définition de la mission}
\subsection{Présentation de l'organisme d'accueil}
Notre projet de fin d'études se déroule au sein de MTD Group, une entreprise spécialisée dans les technologies de l'information et le développement de solutions digitales. MTD Group, dont le nom complet est Media Trade Development, a été créée aux alentours de l'année 2000 et son siège social se trouve à Sfax, en Tunisie, plus précisément à l'adresse suivante : Avenue de Carthage, Immeuble Ribat El Madina, 4\textsuperscript{e} étage, Appartement 405, 3027 Sfax, Tunisie. L'entreprise emploie entre onze et cinquante personnes et s'est imposée au fil des années comme un acteur reconnu du développement logiciel et du digital dans la région.

MTD Group propose une large gamme de services dans le domaine du numérique, notamment la création de sites web, le développement d'applications web et mobiles, le marketing digital et le référencement, ainsi que le développement de solutions ERP pour la gestion d'entreprise, avec sa solution MTD ERP. L'entreprise dispose d'une équipe pluridisciplinaire composée d'ingénieurs, de développeurs, de chefs de projet et de spécialistes du marketing digital, ce qui lui permet de couvrir l'ensemble des besoins de ses clients.

Au‑delà de sa présence à Sfax, MTD Group a également des activités ou des contacts à l'international, notamment en Allemagne à Stuttgart et en Libye à Tripoli, ce qui témoigne de son ouverture sur les marchés étrangers. C'est dans ce cadre que se place notre projet de fin d'études, qui est suivi par un responsable de l'entreprise.

\subsection{Présentation de l'application}
Le projet consiste à créer et développer EduLearn, une application web de gestion de l'apprentissage, que l'on appelle aussi LMS. EduLearn est une plateforme conçue pour un large public, et pas seulement pour MTD Group. Elle pourra être utilisée aussi bien dans des entreprises que dans des écoles ou des universités qui veulent mettre leurs formations en ligne et offrir une expérience d'apprentissage moderne et interactive.

EduLearn se veut une solution facile à utiliser, simple et rapide. C'est une application web adaptable, ce qui signifie qu'elle s'ajuste automatiquement à la taille de l'écran de l'utilisateur, que ce soit un ordinateur de bureau, une tablette ou un smartphone. Il n'y a rien à installer, l'accès se fait directement par un navigateur web comme Chrome, Firefox ou Safari.

Les fonctions principales d'EduLearn répondent à tous les besoins essentiels d'une plateforme d'apprentissage en ligne. Elles comprennent la gestion des utilisateurs avec trois profils différents, la gestion des cours et des leçons, la création et la participation à des quiz, le suivi précis de la progression des apprenants, ainsi que des tableaux de bord pour piloter les activités. L'application sera développée avec des technologies web courantes, en utilisant une architecture MVC qui sépare clairement la présentation, la logique de fonctionnement et l'accès aux données.

\subsection{Objectifs à atteindre}
Les objectifs du projet sont nombreux et concernent à la fois le fonctionnement et la technique.

\begin{itemize}
    \item Rassembler tous les contenus pédagogiques sur une seule plateforme. Il s'agit de mettre fin à l'éparpillement des ressources que l'on rencontre dans beaucoup d'organisations, où les cours sont dispersés sur différents supports et difficiles à trouver.
    \item Automatiser les évaluations. La correction des quiz sera entièrement automatique, ce qui permettra de diminuer fortement le travail administratif des enseignants et des formateurs. Les résultats seront donnés instantanément aux apprenants dès qu'ils auront fini de répondre.
    \item Rendre la plateforme accessible en permanence. Elle devra être disponible vingt‑quatre heures sur vingt‑quatre et sept jours sur sept, pour que les apprenants puissent accéder aux cours à tout moment et de n'importe où, à condition d'avoir une connexion internet.
    \item Suivre la progression de chaque apprenant de façon personnalisée. Des tableaux de bord clairs et des indicateurs utiles permettront aux apprenants de voir où ils en sont et aux enseignants de repérer rapidement ceux qui ont des difficultés.
    \item Simplifier la gestion des utilisateurs. L'administrateur aura une interface spéciale pour créer, modifier, activer ou désactiver les comptes, et pour donner les bons rôles à chaque utilisateur.
    \item Garantir la sécurité. Une gestion précise des droits et une connexion sécurisée protègent les données.
\end{itemize}

\section{Contexte du sujet}
\subsection{Contexte général de la formation en ligne}
La formation en ligne se développe énormément depuis plusieurs années, et cette tendance s'est encore accélérée avec la crise sanitaire récente. Les entreprises et les établissements d'enseignement ont compris qu'il est important d'avoir des outils numériques performants pour assurer la continuité des cours et proposer des formations de qualité, sans être gênés par les distances ou les horaires.

Les plateformes de gestion de l'apprentissage sont devenues des solutions indispensables. Elles permettent d'organiser les parcours de formation, de diffuser des contenus variés, de mettre en place des évaluations et de suivre les progrès des apprenants. Le marché des LMS est en forte croissance, avec une offre très diverse allant des solutions libres aux plateformes commerciales complexes.

\subsection{Les acteurs du système}
Le système EduLearn est prévu pour trois types d'utilisateurs qui ont des rôles bien différents.

\begin{itemize}
    \item L'administrateur est responsable de la gestion générale de la plateforme. Il crée et gère les comptes des utilisateurs, active ou désactive les accès, et regarde les statistiques globales pour vérifier que tout fonctionne bien. Il a une vue d'ensemble sur toutes les activités.
    \item L'enseignant, que l'on appelle aussi formateur, est celui qui crée et gère les contenus pédagogiques. Il prépare les cours, organise les leçons, conçoit les quiz et suit le travail des apprenants. Il peut consulter les résultats et les progrès pour adapter son enseignement.
    \item L'étudiant, ou apprenant, est celui qui utilise la plateforme pour se former. Il regarde les cours auxquels il est inscrit, participe aux quiz, suit sa progression personnelle et consulte l'historique de ses résultats.
\end{itemize}

\section{Analyse d'applications existantes}
Avant de créer notre propre solution, il est nécessaire d'étudier les plateformes déjà disponibles sur le marché. Cette comparaison nous aidera à voir ce qui est bien fait et ce qu'il faut éviter.

\subsection{Moodle}
Moodle est sans aucun doute la plateforme d'apprentissage libre la plus utilisée dans le monde. Créée en 2002 par Martin Dougiamas, elle équipe aujourd'hui plus de deux cent mille institutions dans deux cent quarante‑deux pays. Cette utilisation très large montre qu'elle est fiable et qu'elle a fait ses preuves.

Du point de vue des fonctions, Moodle offre une richesse exceptionnelle. La plateforme permet une gestion très précise des cours, avec de nombreuses possibilités pour organiser et paramétrer les contenus. La gestion des utilisateurs est complète, avec la possibilité de donner différents rôles et de créer des groupes. Les outils d'évaluation sont variés, avec différents types de quiz, des devoirs et des ateliers où les apprenants s'évaluent entre eux. Les fonctions de communication sont aussi très développées, avec des forums, une messagerie et un système de notifications. Le suivi des apprenants se fait grâce à des rapports détaillés et des carnets de notes que l'on peut personnaliser. Enfin, un tableau de bord complet montre des statistiques et des indicateurs utiles.

Cependant, Moodle a plusieurs limites importantes. Son interface est souvent considérée comme complexe et peu facile à prendre en main, ce qui fait qu'il faut beaucoup de temps pour apprendre à s'en servir. L'installation technique est lourde, elle demande des compétences avancées et un serveur spécial pour l'installation et la maintenance. Le grand nombre de fonctions, si c'est un avantage pour les utilisateurs expérimentés, peut perdre ceux qui n'ont besoin que des fonctions de base. Enfin, les mises à jour fréquentes demandent une attention constante de la part des administrateurs.

\subsection{Google Classroom}
Google Classroom est une solution de gestion de l'apprentissage créée par Google et lancée en 2014. Intégrée à l'ensemble Google Workspace for Education, elle est aujourd'hui utilisée par plus de cent cinquante millions d'étudiants et d'enseignants dans le monde, ce qui en fait l'une des plateformes les plus connues, surtout dans l'enseignement primaire et secondaire.

Les points forts de Google Classroom sont sa simplicité d'utilisation et son intégration facile avec les outils Google. On peut créer des classes rapidement et l'inscription des étudiants est simple grâce à des codes. Donner des devoirs est facilité par l'intégration avec Google Drive, qui permet de joindre facilement des documents. Les étudiants peuvent rendre leurs devoirs en ligne et les enseignants peuvent les annoter directement. La communication se fait par un fil d'actualité et une messagerie intégrée.

Malgré ces avantages, Google Classroom a des limites importantes. La dépendance à l'univers Google est totale, ce qui pose problème pour les organisations qui veulent une solution indépendante. Les fonctions de quiz sont assez simples, elles se limitent à des questions à choix multiples avec peu de possibilités de réglage. Les statistiques et rapports sont pauvres, ils donnent peu d'indicateurs pour suivre finement la progression. Enfin, on ne peut rien personnaliser, ni l'interface ni les fonctions.

\subsection{Synthèse comparative}
L'étude de ces deux plateformes montre des profils très différents. Moodle se distingue par sa richesse fonctionnelle exceptionnelle et sa grande souplesse, mais cela demande une utilisation complexe et une technique lourde. Google Classroom séduit par sa simplicité d'utilisation pour les enseignants et les étudiants, mais il a un rôle administrateur limité et une dépendance gênante à Google.

Ces observations nous amènent à placer notre solution EduLearn comme une alternative équilibrée, qui combine les points forts des deux plateformes tout en évitant leurs défauts. Nous visons une simplicité d'utilisation comparable à celle de Google Classroom, une richesse fonctionnelle qui couvre les besoins essentiels comme Moodle, et une indépendance totale par rapport aux fournisseurs extérieurs.

\section{Problématique et solution proposée}
\subsection{Problématique}
Au terme de l'analyse du contexte et des solutions existantes, nous posons la question suivante : comment concevoir et développer une plateforme d'apprentissage en ligne légère, facile à utiliser et performante qui réponde aux besoins des organisations pour gérer les formations et les évaluations, en offrant des interfaces adaptées aux trois profils d'utilisateurs ?

Cette question en soulève plusieurs autres. Quelles fonctions sont vraiment indispensables pour chaque type d'utilisateur ? Comment créer des interfaces à la fois simples et efficaces, adaptées aux besoins particuliers de chacun ? Comment garantir des performances optimales tout en gardant une légèreté technique qui facilite l'installation ? Comment assurer la sécurité des données et une gestion précise des droits selon le rôle ?

\subsection{Solution proposée}
Pour répondre à cette question, nous proposons EduLearn, une plateforme d'apprentissage en ligne générale, conçue pour être utilisée aussi bien en entreprise que dans des établissements d'enseignement. EduLearn a plusieurs qualités essentielles.

\begin{itemize}
    \item Une gestion claire des trois profils d'utilisateurs. Chaque utilisateur a une interface adaptée à ses besoins et des fonctions spécifiques.
    \item La simplicité d'utilisation. L'interface est simple et intuitive, faite pour ne demander aucune formation avant de s'en servir.
    \item La pertinence des fonctions. Nous avons choisi un ensemble de fonctions qui couvre l'essentiel des besoins, sans rien de trop.
    \item La performance. L'architecture est optimisée pour garantir des réponses rapides, même avec beaucoup d'utilisateurs en même temps.
    \item La sécurité. Une gestion précise des droits et une connexion sécurisée protègent les données.
    \item L'autonomie. La solution ne dépend d'aucun fournisseur extérieur, elle est facile à installer et à entretenir.
\end{itemize}

\subsection{Fonctionnalités clés}
Pour l'administrateur, les fonctions principales comprennent la gestion complète des utilisateurs avec l'ajout, la modification, la suppression, l'activation et la désactivation des comptes, la consultation des statistiques générales et la surveillance de la plateforme.

Pour l'enseignant, les fonctions incluent la création, la modification et la suppression de cours, l'ajout de leçons avec différents types de contenu, la publication ou le masquage des cours, la création de quiz liés à un cours, l'ajout de questions, le choix de la note et de la durée, la correction automatique des quiz, la consultation des résultats des étudiants et le suivi de la progression.

Pour l'étudiant, les fonctions comprennent l'inscription et la connexion sécurisées, la consultation des cours disponibles, la participation aux quiz, l'envoi des réponses, la consultation des résultats, le suivi de sa progression personnelle et l'accès à l'historique de ses résultats.

\chapter{Conclusion}
Ce premier chapitre a posé les bases de notre projet de fin d'études. Nous avons présenté l'entreprise qui nous accueille, MTD Group, en précisant son activité et son rôle dans l'encadrement de notre travail. Nous avons défini les grandes lignes de l'application EduLearn, une plateforme générale de gestion de l'apprentissage destinée à un large public. Nous avons identifié les trois profils d'utilisateurs du système et leurs rôles respectifs.

L'étude des applications existantes a montré les forces et les faiblesses de Moodle et de Google Classroom, deux solutions importantes du marché. Moodle impressionne par sa richesse fonctionnelle mais il est trop complexe pour beaucoup d'utilisateurs. Google Classroom séduit par sa simplicité mais il a des limites importantes au niveau des fonctions et de l'indépendance.

Ces observations justifient notre proposition de développer EduLearn, une solution équilibrée qui allie simplicité d'utilisation, fonctions essentielles et indépendance technique. Le prochain chapitre sera consacré à l'étude détaillée des besoins du système.

% -------------------------------------------------
% Bibliographie
\printbibliography

\end{document}
